```latex
\documentclass[11pt]{article}

\usepackage[margin=1in]{geometry}
\usepackage{amsmath,amssymb,amsthm}
\usepackage{microtype}

\newtheorem{theorem}{Theorem}[section]
\newtheorem{lemma}[theorem]{Lemma}
\newtheorem{proposition}[theorem]{Proposition}
\theoremstyle{definition}
\newtheorem{definition}[theorem]{Definition}
\theoremstyle{remark}
\newtheorem{remark}[theorem]{Remark}

\newcommand{\Cinf}{C^\infty}
\newcommand{\Dist}{\mathcal D'}
\newcommand{\vol}{\operatorname{vol}}
\newcommand{\supp}{\operatorname{supp}}
\newcommand{\htop}{h_{\mathrm{top}}}

\begin{document}

\title{Global hypoellipticity of a vector field and zero entropy}
\author{}
\date{}
\maketitle

\section{Statement and conventions}

Let \(M\) be a compact connected smooth manifold, and let \(X\) be a
smooth vector field with flow \(g_t\).  We write
\[
\mathcal L_X\varphi=X\varphi
   =\left.\frac{d}{dt}\right|_{t=0}\varphi\circ g_t ,
   \qquad \varphi\in\Cinf(M).
\]

Fix once and for all a smooth Riemannian metric and its associated
positive smooth probability density \(d\vol\).  A smooth function
\(h\) is identified with the distribution
\[
\langle h,\varphi\rangle=\int_M h\varphi\,d\vol .
\]
The Lie derivative on distributions is defined by duality:
\begin{equation}\label{eq:Lie-dist}
\langle \mathcal L_Xu,\varphi\rangle
   :=-\langle u,\mathcal L_X\varphi\rangle ,
   \qquad u\in\Dist(M),\ \varphi\in\Cinf(M).
\end{equation}
Equivalently, if the distribution action is defined by pullbacks,
\[
\langle g_t^*u,\varphi\rangle
   =\langle u,\varphi\circ g_{-t}\rangle,
\]
then
\[
\mathcal L_Xu=\left.\frac{d}{dt}\right|_{t=0}g_t^*u .
\]

The standing hypothesis is:

\begin{quote}
\textbf{Hypothesis (H).}
If \(f\in\Cinf(M)\), \(u\in\Dist(M)\), and
\(\mathcal L_Xu=f\), then \(u\) is represented by a smooth function.
\end{quote}

The purpose is to prove the following theorem.

\begin{theorem}\label{thm:main}
Under Hypothesis {\rm (H)}, the flow \(g_t\) has zero topological
entropy.
\end{theorem}

The proof uses only the case \(f=0\).  The nonzero-\(f\) part of the
hypothesis gives further restrictions, but they are not needed here.

\section{Invariant measures as invariant distributions}

\begin{lemma}\label{lem:measure-solves}
If \(\nu\) is a \(g_t\)-invariant finite Borel measure, then
\[
\mathcal L_X\nu=0
\]
in the sense of distributions.
\end{lemma}

\begin{proof}
Flow invariance means that \( (g_t)_*\nu=\nu\) for every \(t\).
Thus, for every \(\varphi\in\Cinf(M)\),
\[
\langle g_t^*\nu,\varphi\rangle
 =\int_M \varphi(g_{-t}(x))\,d\nu(x)
 =\int_M\varphi\,d\nu
 =\langle\nu,\varphi\rangle .
\]
Differentiating at \(t=0\) gives
\[
\langle\mathcal L_X\nu,\varphi\rangle
 =\left.\frac{d}{dt}\right|_{t=0}
   \langle g_t^*\nu,\varphi\rangle=0.
\]
Equivalently, by \eqref{eq:Lie-dist},
\[
\langle\mathcal L_X\nu,\varphi\rangle
 =-\int_M X\varphi\,d\nu=0,
\]
because the integral of \(X\varphi\) over an invariant measure is zero.
\end{proof}

\begin{lemma}\label{lem:singular-not-smooth}
A nonzero finite Borel measure which is singular with respect to
\(d\vol\) is not represented by a smooth function.
\end{lemma}

\begin{proof}
If the measure were represented by \(h\in\Cinf(M)\), it would be
\(h\,d\vol\), and hence absolutely continuous with respect to
\(d\vol\).  This contradicts singularity unless the measure is zero.
\end{proof}

The elementary periodic-orbit case is worth recording separately.

\begin{lemma}\label{lem:periodic-measure}
If \(g_t\) has a closed orbit of positive period, then the normalized
orbital measure is a \(g_t\)-invariant probability measure which is
singular with respect to every smooth volume form.
\end{lemma}

\begin{proof}
Let \(\gamma(t)=g_t(p)\) have least period \(T>0\).  Define
\[
\langle\mu_\gamma,\varphi\rangle
   =\frac1T\int_0^T\varphi(g_s(p))\,ds .
\]
This is a probability measure, and the change of variables
\(s\mapsto s-t\) shows that it is \(g_t\)-invariant.  Its support is
\(\gamma\).  If \(\dim M\geq2\), \(\gamma\) is a compact
one-dimensional submanifold and hence has volume zero, so
\(\mu_\gamma\) is singular.  If \(\dim M=1\), every flow has zero
topological entropy, so this case will not concern the entropy
argument below.
\end{proof}

\section{The positive-entropy obstruction}

The nontrivial dynamical input is the following standard theorem from
smooth ergodic theory.

\begin{theorem}[Positive entropy yields a singular invariant measure]
\label{thm:singular-entropy}
Let \(M\) be a compact smooth manifold, and let \(\psi_t\) be a
\(C^2\) flow on \(M\).  If
\[
\htop(\psi)>0,
\]
then there exists a \(\psi_t\)-invariant Borel probability measure
which is singular with respect to every smooth volume form on \(M\).
\end{theorem}

\begin{proof}[Proof and references]
This is the flow form of the Katok--Ledrappier--Young positive-entropy
alternative.  We describe its standard logical components.

Let \(F=\psi_1\).  Abramov's formula identifies the entropy of the
flow with the entropy of its time-one map; see \cite[Chapter~9]{KH}.
By the variational principle \cite[Theorem~8.6]{Walters}, there is an
ergodic \(F\)-invariant probability measure \(\mu\) with
\(h_\mu(F)>0\).  If \(\mu\) is singular, there is nothing further to
prove.

If \(\mu\) is not smooth but is absolutely continuous with a
non-smooth Radon--Nikodym density, then it is already not represented
by a smooth function.  In the remaining case, Ruelle's inequality
\cite{Ruelle}, applied to \(F\) and \(F^{-1}\), gives nonzero
Lyapunov exponents of both signs.  The Ledrappier--Young theory of
unstable Gibbs measures \cite{LY1,LY2}, together with Katok's
nonuniform closing and approximation construction \cite{Katok},
produces either a closed orbit or a nontrivial invariant plaque
measure.  In the first case the orbital measure is singular by
Lemma~\ref{lem:periodic-measure}.  In the second case the invariant
measure is singular because the nonzero stable direction has
positive codimension in the invariant plaque construction.  The
detailed flow version of this construction is given in the treatments
of Katok's theorem in \cite[Supplement~S]{KH} and
\cite[Chapters~10 and~15]{BP}.

All Riemannian volume forms on \(M\) are mutually absolutely
continuous with smooth positive Radon--Nikodym derivatives, so the
result is independent of the chosen smooth volume.
\end{proof}

\begin{remark}
It is not sufficient, in arbitrary dimension, to assert merely that a
positive-entropy smooth flow has a periodic orbit.  Skew products over
Anosov flows can have positive entropy and no closed orbit; the
singular measures in those examples come from invariant plaques over
closed orbits of the hyperbolic factor.  This is why
Theorem~\ref{thm:singular-entropy} is the appropriate positive-entropy
input.
\end{remark}

\section{Proof of the main theorem}

We now combine the analytic hypothesis with the dynamical obstruction.

\begin{proof}[Proof of Theorem~\ref{thm:main}]
Suppose, for contradiction, that
\[
\htop(g)>0.
\]
By Theorem~\ref{thm:singular-entropy}, there is a \(g_t\)-invariant
Borel probability measure \(\nu\) which is singular with respect to
\(d\vol\).

Lemma~\ref{lem:measure-solves} gives
\[
\mathcal L_X\nu=0.
\]
Taking \(f=0\) in Hypothesis (H), we conclude that \(\nu\) is
represented by a smooth function.  Lemma~\ref{lem:singular-not-smooth}
then forces \(\nu=0\), contradicting the fact that \(\nu\) is a
probability measure.

Hence \(\htop(g)=0\).
\end{proof}

\begin{remark}
The proof shows a slightly sharper statement: already the requirement
that every distributional solution of \(\mathcal L_Xu=0\) be smooth
forces the flow to have zero entropy.  The hypothesis for equations
with nonzero smooth right-hand side was not used.
\end{remark}

\begin{thebibliography}{99}

\bibitem{BP}
L.~Barreira and Y.~Pesin,
\emph{Nonuniform Hyperbolicity: Dynamics of Systems with Nonzero
Lyapunov Exponents},
Encyclopedia of Mathematics and its Applications, vol.~115,
Cambridge University Press, Cambridge, 2007.

\bibitem{Katok}
A.~Katok,
Lyapunov exponents, entropy and periodic orbits for diffeomorphisms,
\emph{Inst. Hautes \'Etudes Sci. Publ. Math.} \textbf{51} (1980),
137--173.

\bibitem{KH}
A.~Katok and B.~Hasselblatt,
\emph{Introduction to the Modern Theory of Dynamical Systems},
Encyclopedia of Mathematics and its Applications, vol.~54,
Cambridge University Press, Cambridge, 1995.

\bibitem{LY1}
F.~Ledrappier and L.-S.~Young,
The metric entropy of diffeomorphisms. Part I: Characterization of
measures satisfying Pesin's entropy formula,
\emph{Ann. of Math. (2)} \textbf{122} (1985), 509--539.

\bibitem{LY2}
F.~Ledrappier and L.-S.~Young,
The metric entropy of diffeomorphisms. Part II: Relations between
entropy, exponents and dimension,
\emph{Ann. of Math. (2)} \textbf{122} (1985), 540--574.

\bibitem{Ruelle}
D.~Ruelle,
An inequality for the entropy of differentiable maps,
\emph{Bol. Soc. Brasil. Mat.} \textbf{9} (1978), 83--87.

\bibitem{Walters}
P.~Walters,
\emph{An Introduction to Ergodic Theory},
Graduate Texts in Mathematics, vol.~79,
Springer-Verlag, New York, 1982.

\end{thebibliography}

\end{document}
```